% ============================================================
% Paper 94 (EN): Hartshorne's Conjectures in Algebraic Geometry
% Author: Michael Fuhrmann
% Project: Specialist Mathematics Project
% Build: 169
% Date: 2026-03-12
% ============================================================
\documentclass[12pt,a4paper]{amsart}

\usepackage{amsmath,amssymb,amsthm,mathtools}
\usepackage{enumitem}
\usepackage{booktabs}
\usepackage{array}
\usepackage[hidelinks]{hyperref}
\usepackage[utf8]{inputenc}
\usepackage[T1]{fontenc}
\usepackage{geometry}
\geometry{margin=2.5cm}

% ---- Theorem Environments -----------------------------------
\newtheorem{theorem}{Theorem}[section]
\newtheorem{lemma}[theorem]{Lemma}
\newtheorem{corollary}[theorem]{Corollary}
\newtheorem{proposition}[theorem]{Proposition}
\theoremstyle{definition}
\newtheorem{definition}[theorem]{Definition}
\newtheorem{example}[theorem]{Example}
\theoremstyle{remark}
\newtheorem{remark}[theorem]{Remark}
\newtheorem{conjecture}[theorem]{Conjecture}

% ---- Custom Operators ---------------------------------------
\DeclareMathOperator{\codim}{codim}
\DeclareMathOperator{\rank}{rank}
\DeclareMathOperator{\Ext}{Ext}
\DeclareMathOperator{\Hom}{Hom}
\DeclareMathOperator{\Pic}{Pic}
\DeclareMathOperator{\Spec}{Spec}
\DeclareMathOperator{\Proj}{Proj}
\newcommand{\PP}{\mathbb{P}}
\newcommand{\ZZ}{\mathbb{Z}}
\newcommand{\CC}{\mathbb{C}}
\newcommand{\RR}{\mathbb{R}}
\newcommand{\OO}{\mathcal{O}}
\newcommand{\FF}{\mathcal{F}}

% ---- Abstract -----------------------------------------------
\begin{abstract}
We survey Hartshorne's conjectures on complete intersections and vector
bundles in projective space, with emphasis on rigorous statements and the
current state of knowledge.  The Quillen--Suslin theorem (1976), which
resolved Serre's conjecture that projective modules over polynomial rings
are free, is presented in detail.  We discuss the Barth--Lefschetz theorem
on cohomological connectivity, the Horrocks--Mumford rank-2 bundle on $\PP^4$,
the splitting conjecture for low-rank vector bundles on $\PP^n$, and the
Hartshorne--Lichtenbaum vanishing theorem.  Hartshorne's codimension-2
conjecture---that every smooth codimension-2 subvariety of $\PP^n$ for
$n \ge 7$ is a complete intersection---remains open.
\end{abstract}

\title{Hartshorne's Conjectures in Algebraic Geometry}

\author{Michael Fuhrmann}
\address{Specialist Mathematics Project, Build 169}
\email{specialist-maths@project.local}
\date{2026-03-12}

\maketitle
\tableofcontents

% =============================================================
\section{Introduction}
% =============================================================

Projective space $\PP^n_\CC$ is the simplest compact algebraic variety, yet
understanding which subvarieties can live inside it remains deep.  Among the
most elegant questions: which smooth subvarieties are complete intersections?

A smooth subvariety $Y \subseteq \PP^n$ of codimension $c$ is a
\emph{complete intersection} if its ideal $I(Y)$ is generated by exactly
$c$ homogeneous polynomials.  Complete intersections are well understood by
commutative algebra; arbitrary subvarieties are far more mysterious.

In his 1974 paper \cite{Hartshorne1974} and book \cite{Hartshorne1977},
Hartshorne formulated several conjectures suggesting that in high-dimensional
projective space, smooth subvarieties of small codimension must be complete
intersections.  These conjectures encode a deep "rigidity" of $\PP^n$.

Concurrently, Serre's conjecture that projective modules over polynomial rings
$k[x_1,\ldots,x_n]$ are free---equivalent to asking whether every vector
bundle on affine $n$-space is trivial---was proved independently by Quillen
\cite{Quillen1976} and Suslin \cite{Suslin1976} in 1976.

This paper presents rigorous accounts of all these results, distinguishing
carefully between what is proved and what remains conjectural.

% =============================================================
\section{Complete Intersections in Projective Space}
% =============================================================

\begin{definition}[Complete Intersection]
A closed subscheme $Y \subseteq \PP^n_k$ of codimension $c$ is a
\emph{(scheme-theoretic) complete intersection} if the ideal sheaf
$\mathcal{I}_Y$ is generated by exactly $c$ elements, i.e.,
$\mathcal{I}_Y = (f_1,\ldots,f_c)$ for homogeneous polynomials
$f_i \in k[x_0,\ldots,x_n]$ of degrees $d_1 \le \cdots \le d_c$.
\end{definition}

\begin{example}
\begin{enumerate}[label=(\roman*)]
  \item Hypersurfaces in $\PP^n$ (codimension 1) are always complete
        intersections.
  \item Smooth curves in $\PP^3$ (codimension 2): the twisted cubic
        $C = \{[s^3:s^2t:st^2:t^3]\} \subseteq \PP^3$ has degree 3 and
        arithmetic genus 0.  It is \emph{not} a complete intersection;
        its ideal requires 3 generators (the $2\times 2$ minors of a
        $2\times 3$ matrix).
  \item The Veronese surface $V \subseteq \PP^5$ (codimension 3) is not a
        complete intersection.
\end{enumerate}
\end{example}

\begin{theorem}[Graded Nakayama Lemma and Codimension Bounds]
Over an algebraically closed field, a smooth complete intersection $Y$ of
codimension $c$ in $\PP^n$ satisfies $c \le n$.  The canonical class is
$K_Y = \OO_Y\!\bigl(\sum d_i - (n+1)\bigr)$.
\end{theorem}

% =============================================================
\section{The Barth--Lefschetz Theorem}
% =============================================================

\begin{theorem}[Barth, 1970 \cite{Barth1970}]
\label{thm:barth}
Let $Y \subseteq \PP^n_\CC$ be a smooth subvariety of codimension $c$.
Then the restriction map
\[
  H^k(\PP^n, \ZZ) \longrightarrow H^k(Y, \ZZ)
\]
is an isomorphism for $k < \dim Y - c + 1 = n - 2c + 1$ and injective for
$k = n - 2c + 1$.
\end{theorem}

\begin{remark}
For a complete intersection $Y$ of codimension $c$ in $\PP^n$, the
Lefschetz hyperplane theorem gives isomorphisms in an even wider range:
$k \le \dim Y - 1 = n - c - 1$.  Barth's theorem is the analogue without
the complete intersection hypothesis, but at the cost of a more restrictive
bound depending on $c$.
\end{remark}

\begin{corollary}
If $Y \subseteq \PP^n$ is smooth with $\codim Y = c$ and $n \ge 2c$, then
$\pi_1(Y) = 1$ and $\Pic(Y) \cong \ZZ$ (generated by $\OO_Y(1)$).
\end{corollary}

The Barth--Lefschetz theorem is particularly powerful because it shows that
subvarieties of small codimension in large $\PP^n$ are topologically
"simple."  Complete intersections would be even simpler; Hartshorne's
conjecture asserts that the converse should hold.

% =============================================================
\section{Hartshorne's Codimension-2 Conjecture}
% =============================================================

\begin{conjecture}[Hartshorne, 1974 \cite{Hartshorne1974}, Open]
\label{conj:hartshorne}
Let $Y \subseteq \PP^n_\CC$ be a smooth subvariety of codimension 2.  If
$n \ge 7$, then $Y$ is a complete intersection.
\end{conjecture}

\begin{remark}
The bound $n \ge 7$ is essential:
\begin{itemize}
  \item For $n = 3$: smooth codimension-2 subvarieties are curves in $\PP^3$,
        and many non-complete-intersection curves exist (e.g., twisted cubic).
  \item For $n = 4$: the Horrocks--Mumford surface exists as a non-CI
        surface in $\PP^4$ (abelian surface of degree 10).
  \item For $n = 5, 6$: the conjecture is also open (and some evidence
        suggests it may fail for $n = 5$ or 6 as well).
  \item For $n \ge 7$: Hartshorne's conjecture (open, no counterexamples known).
\end{itemize}
\end{remark}

\begin{theorem}[Partial Evidence, Hartshorne--Ogus \cite{HartshorneOgus1974}]
If $Y \subseteq \PP^n$ is smooth of codimension 2 and $n \ge 2\dim Y + 1$,
then the normal bundle $N_{Y/\PP^n}$ splits as a direct sum of line bundles.
In this range, $Y$ has the cohomological properties of a complete intersection.
\end{theorem}

The connection to vector bundles is key: a smooth codimension-2 subvariety
$Y \subseteq \PP^n$ is a complete intersection if and only if the ideal
sheaf $\mathcal{I}_Y$ has a free resolution of length 2 (by the
Hilbert syzygy theorem in this range), which corresponds to the existence of
a rank-2 vector bundle fitting into the \emph{Serre correspondence}.

% =============================================================
\section{The Serre Correspondence and Vector Bundles}
% =============================================================

\begin{theorem}[Serre Correspondence]
Let $Y \subseteq \PP^n$ be a locally Cohen--Macaulay codimension-2 subvariety.
Then there exists a rank-2 vector bundle $\mathcal{E}$ on $\PP^n$ and a
section $s \in H^0(\PP^n, \mathcal{E}(t))$ for some $t$, such that
$Y = Z(s)$ (the zero locus of $s$).

Conversely, the zero locus of a general section of a rank-2 bundle on
$\PP^n$ (for $n \ge 3$) is a smooth codimension-2 subvariety.
\end{theorem}

\begin{corollary}
Every smooth codimension-2 subvariety of $\PP^n$ arises as the zero locus
of a section of a rank-2 vector bundle.  Hartshorne's conjecture is therefore
equivalent to: every rank-2 bundle on $\PP^n$ ($n \ge 7$) is a direct sum
of line bundles $\OO(a) \oplus \OO(b)$.
\end{corollary}

% =============================================================
\section{The Splitting Conjecture for Vector Bundles}
% =============================================================

\begin{conjecture}[Splitting Conjecture, Open]
\label{conj:splitting}
Every algebraic vector bundle of rank $r \le \lfloor n/2 \rfloor$ on
$\PP^n_\CC$ splits as a direct sum of line bundles.
\end{conjecture}

\begin{theorem}[Horrocks Criterion \cite{Horrocks1964}]
A vector bundle $\mathcal{E}$ on $\PP^n$ splits as a direct sum of line bundles
if and only if $H^i(\PP^n, \mathcal{E}(k)) = 0$ for all $k \in \ZZ$ and
$1 \le i \le n-1$.
\end{theorem}

\begin{theorem}[Rank-1 bundles on $\PP^n$]
Every rank-1 algebraic vector bundle on $\PP^n_\CC$ is isomorphic to
$\OO(d)$ for a unique $d \in \ZZ$.  (This follows from $\Pic(\PP^n) \cong \ZZ$.)
\end{theorem}

\begin{theorem}[Serre--Swan, see \cite{Hartshorne1977}]
For $n \ge 3$, every rank-2 vector bundle on $\PP^n_\CC$ that restricts to
a split bundle on every line is itself split (Horrocks's result by Horrocks
criterion and Leray spectral sequence).
\end{theorem}

The situation depends critically on the rank relative to $n$:

\begin{center}
\begin{tabular}{clll}
\toprule
$n$ & rank $r$ & Status & Reference \\
\midrule
$\PP^1$ & any & Always split & Classical \\
$\PP^2$ & any & Non-split exist ($r \ge 2$) & Tautological bundle \\
$\PP^3$ & $r=2$ & Non-split exist & Tautological, null-correlation \\
$\PP^4$ & $r=2$ & Horrocks--Mumford exists & \cite{HorrocksMumford1973} \\
$\PP^5$ & $r=2$ & Open & --- \\
$\PP^n$, $n \ge 6$ & $r=2$ & Splitting conjecture (open) & \cite{Hartshorne1974} \\
\bottomrule
\end{tabular}
\end{center}

% =============================================================
\section{The Horrocks--Mumford Bundle}
% =============================================================

\begin{theorem}[Horrocks--Mumford, 1973 \cite{HorrocksMumford1973}]
There exists a rank-2 indecomposable vector bundle $\mathcal{F}_\mathrm{HM}$
on $\PP^4_\CC$, the \emph{Horrocks--Mumford bundle}, with Chern classes
$c_1 = 5$, $c_2 = 10$.  The zero loci of its sections are abelian surfaces
of degree 10 in $\PP^4$.
\end{theorem}

\begin{remark}
The Horrocks--Mumford bundle is the unique known indecomposable rank-2 bundle
on $\PP^n$ for $n \ge 4$.  Its existence shows that Hartshorne's splitting
conjecture fails for $n = 4$.  The question for $n \ge 5$ is open.
\end{remark}

The construction of $\mathcal{F}_\mathrm{HM}$ uses the action of the
Heisenberg group $H_5$ on $\PP^4$ and the theory of theta functions.
The zero loci are Horrocks--Mumford abelian surfaces.

% =============================================================
\section{Serre's Conjecture: The Quillen--Suslin Theorem}
% =============================================================

\begin{theorem}[Quillen--Suslin, 1976 \cite{Quillen1976,Suslin1976}]
\label{thm:quillen-suslin}
Let $k$ be a field (or a principal ideal domain).  Every finitely generated
projective module over the polynomial ring $k[x_1,\ldots,x_n]$ is free.
\end{theorem}

\begin{remark}
Serre formulated this as a conjecture in 1955 \cite{Serre1955}.
Geometrically, it says: every algebraic vector bundle on the affine space
$\mathbb{A}^n_k$ is trivial (isomorphic to $\OO^r$ for some $r$).
This is in sharp contrast to the projective case.
\end{remark}

\begin{proof}[Proof sketch (Quillen's approach)]
\textbf{Step 1.}  By a theorem of Horrocks (1964), a projective module $P$
over $k[x_1,\ldots,x_n]$ is free if and only if $P \otimes_{k[x_n]} k[x_n]_{(f)}$
is free for a single monic polynomial $f \in k[x_n]$.

\textbf{Step 2.}  Quillen's \emph{Patching Lemma}: if $M$ is a finitely
presented $A[t]$-module such that $M_\mathfrak{m}$ is extended from
$A_\mathfrak{m}$ for all maximal ideals $\mathfrak{m} \subset A$, then $M$
is extended from $A$.

\textbf{Step 3.}  Induction on $n$: for $n=1$, projective modules over a
PID are free.  The induction uses the Patching Lemma to reduce to the
$n-1$ case. $\square$
\end{proof}

The Quillen--Suslin theorem has become fundamental in algebraic K-theory.
Suslin's proof was more elementary (matrix calculations), while Quillen's
proof introduced powerful new techniques.

% =============================================================
\section{The Hartshorne--Lichtenbaum Vanishing Theorem}
% =============================================================

\begin{theorem}[Hartshorne--Lichtenbaum Vanishing Theorem \cite{Hartshorne1968}]
\label{thm:HLVT}
Let $X$ be a smooth projective variety of dimension $n$ over a field $k$,
and let $Y \subseteq X$ be a closed subscheme.  Then
\[
  H^n_Y(X, \OO_X) = 0
\]
provided that $Y$ does not contain any irreducible component of $X$ and
$\codim_X(Y) \ge 1$.  More precisely, for an ample line bundle $\mathcal{L}$:
\[
  H^i(X, \mathcal{L}^{-1}) = 0 \quad \text{for } i < n.
\]
\end{theorem}

\begin{remark}
This generalises the Kodaira vanishing theorem ($H^i(X, K_X \otimes L) = 0$
for $L$ ample and $i > 0$) and is a cornerstone of modern algebraic geometry.
In characteristic $p > 0$, analogues require additional hypotheses.
\end{remark}

\begin{corollary}
Combined with the Barth--Lefschetz theorem, the HLVT implies strong
restrictions on the topology of subvarieties: a smooth subvariety of
codimension 2 in $\PP^n$ for large $n$ has the same cohomology as a
complete intersection in many degrees.
\end{corollary}

% =============================================================
\section{Algebraic K-Theory and Vector Bundles}
% =============================================================

Serre's conjecture and the splitting conjecture both fit into the framework
of algebraic K-theory.

\begin{definition}[Grothendieck Group]
For a scheme $X$, the \emph{Grothendieck group} $K_0(X)$ is the free abelian
group on isomorphism classes of coherent sheaves on $X$, modulo the relation
$[\mathcal{F}] = [\mathcal{F}'] + [\mathcal{F}'']$ for every short exact
sequence $0 \to \mathcal{F}' \to \mathcal{F} \to \mathcal{F}'' \to 0$.
\end{definition}

\begin{theorem}[Quillen, 1973 \cite{Quillen1973}]
$K_i(\PP^n_k) \cong K_i(k)^{n+1}$ for all $i \ge 0$ (the projective bundle
formula in higher K-theory).
\end{theorem}

For affine space: $K_i(\mathbb{A}^n_k) \cong K_i(k)$ for all $i$ (homotopy
invariance), which is the K-theoretic form of the Quillen--Suslin theorem.

% =============================================================
\section{Current Status and Related Results}
% =============================================================

\begin{theorem}[Zak's Theorem on Tangencies, 1981 \cite{Zak1993}]
Let $Y \subseteq \PP^n$ be a smooth non-degenerate subvariety.  If
$\dim Y > \frac{2}{3}(n-1)$, then $Y$ cannot be tangent to a hyperplane
along a variety of dimension $> 0$.  In particular, Scorza varieties and
Severi varieties are classified.
\end{theorem}

\begin{theorem}[Faltings, 1981 \cite{Faltings1981}]
If $Y \subseteq \PP^n$ is a smooth complete intersection of codimension $c$
and $n > 2c$, then $Y$ is a scheme-theoretic complete intersection (not just
set-theoretic).
\end{theorem}

\begin{theorem}[Evans--Griffith, 1981 \cite{EvansGriffith1981}]
Let $\mathcal{E}$ be a vector bundle of rank $r$ on $\PP^n_k$ (algebraically
closed field).  If $\mathcal{E}$ splits on every $(r-1)$-dimensional projective
subspace, then $\mathcal{E}$ splits.
\end{theorem}

% =============================================================
\section{Open Problems}
% =============================================================

\begin{conjecture}[Hartshorne's Codimension-2 Conjecture, = Conjecture~\ref{conj:hartshorne}, Open]
Every smooth codimension-2 subvariety $Y \subseteq \PP^n_\CC$ for $n \ge 7$
is a complete intersection.
\end{conjecture}

\begin{conjecture}[Splitting Conjecture for Rank-2 Bundles, = Conjecture~\ref{conj:splitting}, Open]
Every rank-2 algebraic vector bundle on $\PP^n_\CC$ for $n \ge 6$
decomposes as $\OO(a) \oplus \OO(b)$.
\end{conjecture}

\begin{conjecture}[Set-theoretic Complete Intersections, Partially Open]
Every smooth subvariety $Y \subseteq \PP^n$ is a \emph{set-theoretic}
complete intersection (i.e., defined as a set by $\codim Y$ equations).
This is known in characteristic $p > 0$ by Kempf (using Frobenius), but
open in characteristic 0 for curves in $\PP^3$.
\end{conjecture}

\begin{conjecture}[Hartshorne's General Conjecture \cite{Hartshorne1974}]
If $Y$ is a smooth subvariety of $\PP^n$ with $\dim Y > \frac{2n}{3}$,
then $Y$ is a complete intersection.
\end{conjecture}

% =============================================================
\section{Conclusion}
% =============================================================

Hartshorne's conjectures have driven a generation of research at the
interface of algebraic geometry, commutative algebra, and topology.  The
Quillen--Suslin theorem resolved the affine case completely; the projective
case via the splitting conjecture and Hartshorne's codimension-2 conjecture
remains substantially open.

The Horrocks--Mumford bundle demonstrates the subtlety of the projective
case: indecomposable rank-2 bundles exist on $\PP^4$, but none are known on
$\PP^n$ for $n \ge 5$.  Whether this reflects a genuine geometric obstruction
or merely a gap in our techniques is unknown.

\begin{thebibliography}{99}

\bibitem{Hartshorne1974}
R.~Hartshorne,
\textit{Varieties of small codimension in projective space},
Bull. Amer. Math. Soc. \textbf{80} (1974), 1017--1032.

\bibitem{Hartshorne1977}
R.~Hartshorne,
\textit{Algebraic Geometry},
Graduate Texts in Mathematics \textbf{52},
Springer, New York, 1977.

\bibitem{Quillen1976}
D.~Quillen,
\textit{Projective modules over polynomial rings},
Invent. Math. \textbf{36} (1976), 167--171.

\bibitem{Suslin1976}
A.~A. Suslin,
\textit{Projective modules over a polynomial ring},
Dokl. Akad. Nauk SSSR \textbf{226} (1976), 549--551
(English transl.: Soviet Math. Dokl. \textbf{17} (1976), 1160--1164).

\bibitem{Barth1970}
W.~Barth,
\textit{Transplanting cohomology classes in complex-projective space},
Amer. J. Math. \textbf{92} (1970), 951--967.

\bibitem{HorrocksMumford1973}
G.~Horrocks and D.~Mumford,
\textit{A rank 2 vector bundle on $\mathbf{P}^4$ with 15{,}000 symmetries},
Topology \textbf{12} (1973), 63--81.

\bibitem{Horrocks1964}
G.~Horrocks,
\textit{Vector bundles on the punctured spectrum of a local ring},
Proc. London Math. Soc. (3) \textbf{14} (1964), 689--713.

\bibitem{Serre1955}
J.-P. Serre,
\textit{Faisceaux algébriques cohérents},
Ann. of Math. (2) \textbf{61} (1955), 197--278.

\bibitem{Hartshorne1968}
R.~Hartshorne,
\textit{Cohomological dimension of algebraic varieties},
Ann. of Math. (2) \textbf{88} (1968), 403--450.

\bibitem{HartshorneOgus1974}
R.~Hartshorne and A.~Ogus,
\textit{On the factoriality of local rings of small embedding codimension},
Comm. Algebra \textbf{1} (1974), 415--437.

\bibitem{Quillen1973}
D.~Quillen,
\textit{Higher algebraic K-theory I},
in: \textit{Algebraic K-theory I}, Lecture Notes in Math. \textbf{341},
Springer, 1973, 85--147.

\bibitem{Faltings1981}
G.~Faltings,
\textit{Ein Kriterium für vollständige Durchschnitte},
Invent. Math. \textbf{62} (1981), 393--401.

\bibitem{EvansGriffith1981}
E.~G. Evans and P.~Griffith,
\textit{The syzygy problem},
Ann. of Math. (2) \textbf{114} (1981), 323--333.

\bibitem{Zak1993}
F.~L. Zak,
\textit{Tangents and Secants of Algebraic Varieties},
Translations of Math. Monographs \textbf{127},
American Mathematical Society, 1993.

\end{thebibliography}

\end{document}
